%% P15 --- Functions of several variables and the gradient
%%
%% Every number in this program that the reader cannot do in their head comes
%% from code/p15_gradient.py and is pulled in with \val{}. The console listing
%% is a committed transcript written by that script.
%%
%% THE THESIS: the gradient's DIRECTION is derived rather than chosen. The
%% directional derivative is a dot product, so Program P05's cosine settles
%% which direction is steepest and which are flat, and "steepest" and
%% "perpendicular to the contour" turn out to be two readings of one equation.
%%
%% WHAT THIS PROGRAM IS OWED, read out of the files rather than remembered:
%%   F11  derives the one-dimensional walk IN FULL -- downhill is against the
%%        sign of the derivative, and x <- x - eta f'(x) with the minus sign
%%        called "the entire steering mechanism". So this program MAY NOT
%%        re-derive the sign. What it owes back is the same recurrence per
%%        eigendirection, which the script gates.
%%   F12  says d/dx names WHICH letter is the variable and that this "will
%%        matter enormously in Program P15, where there is more than one".
%%   P05  gives a . b = |a| |b| cos(theta). Section 3 is that identity with
%%        one of the two vectors being the gradient. Nothing else is needed.
%%   P10  collects the shape of the basin -- eigenvalues 20 and 1, an ellipse
%%        4.47 times longer than wide -- and says explicitly that it does not
%%        spend it. Section 5 spends it, on P10's own committed numbers.
%%
%% WHAT IT LEAVES ALONE, checked against tools/programs.json:
%%   Jacobians and reverse mode                                  -> P16
%%   the Hessian, curvature, and THE LARGEST STABLE STEP          -> P17
%%   matrix calculus, the softmax-cross-entropy gradient          -> P18
%%   convexity                                                    -> P19
%%   momentum, which FIXES the zig-zag this program produces      -> P20
%% P20's brief says in as many words that it fixes "the zig-zag from P15", so
%% producing it and leaving it unfixed is the contract rather than a gap.
%%
%% THE WORD COLLIDES, and section 2 says so. F06 and F11 use "gradient" for
%% the SLOPE OF A LINE, which is standard British usage and is what those
%% programs needed; here it is a vector. Same treatment P07 gave "dimension"
%% doing three jobs.
%%
%% Twin: programs/pl/P15-gradient.tex. Frame for frame, macro for macro,
%% number for number; tools/parity.py compares them.

\program{Functions of several variables and the gradient}
\label{prog:P15}
\index{gradient}
\index{derivative!partial}

\begin{outcomes}
\outcome{1--7}{take a partial derivative by holding every other input still,
  and say which question it answers}
\outcome{8--13}{collect the partials into a gradient, and keep the vector
  apart from the slope of a line}
\outcome{14--19}{write the rate of change along any direction as a dot
  product with the gradient}
\outcome{20--28}{say why the gradient is the steepest direction and why it
  meets a contour at right angles, from that one dot product}
\outcome{29--38}{say what the walk does in a long thin valley, and why}
\end{outcomes}

\begin{quiz}
\item $f(x, y) = x^{2}y + 3y$. What is $\partial f/\partial x$?
  \teachesat{3--4}
  \answerto{$2xy$. Treat $y$ as a constant and differentiate as
  Program~\ref{prog:F11} taught \dash{} the $3y$ term has no $x$ in it, so it
  contributes nothing.}
\item Which direction, out of all of them, does $f$ increase fastest in?
  \teachesat{20--21}
  \answerto{The gradient's own direction. Not because it is defined that way
  \dash{} because the rate along a direction is the gradient dotted with it,
  and a dot product is largest when the two point the same way.}
\item Gradient descent moves in which direction? \teachesat{29--30}
  \answerto{The \emph{negative} gradient. The gradient points the way $f$
  increases fastest, so the step goes against it, and the sign is the
  commonest error in a hand-written training loop.}
\item You are on a contour line of $f$ and walk along it. What is the rate of
  change of $f$ along your path? \teachesat{23--24}
  \answerto{Zero, by the definition of a contour. And since that rate is the
  gradient dotted with your direction, the gradient must be at right angles
  to the contour.}
\item In a long thin valley, does the negative gradient point at the lowest
  point? \teachesat{31--32}
  \answerto{No, and that is the whole of section 5. It points across the
  valley, at right angles to the contour you are standing on, which is not
  the same direction as the one home unless the bowl is round.}
\end{quiz}

% ============================================================
\section{One input at a time}
\label{sec:P15-partials}
\index{derivative!partial}

\begin{fr}
Everything Program~\ref{prog:F11} and Program~\ref{prog:F12} did was about a
function of one input. A loss is not. It is a function of every parameter in
the model at once, which for anything worth training is a number with nine or
ten digits in it.

That sounds like a different subject and it is not. Program~\ref{prog:F12}
left a clue: it said $\frac{\mathrm{d}}{\mathrm{d}x}$ names \emph{which}
letter is the variable, and that naming it would matter here.

Given $f(x, y)$, and only the machinery of one-variable calculus, how would
you get a derivative out of it at all?
\dotline
\nextframe
\end{fr}

\begin{fr}
\begin{ansblock}
Hold one of them still. Fix $y$ at some value and $f$ becomes a function of
$x$ alone, which is a thing you can already differentiate.
\end{ansblock}

That is the whole idea and there is nothing more to it. It is written
$\frac{\partial f}{\partial x}$, with a rounded $\partial$ rather than a
straight $\mathrm{d}$, and read \enquote{partial dee $f$ by dee $x$}. The
rounded letter is a reminder that something else is being held still.

\begin{note}
The notation carries no new mathematics. $\frac{\partial f}{\partial x}$ is
an ordinary derivative of an ordinary one-variable function \dash{} the
function you get from $f$ by freezing every input except $x$. Every rule
Program~\ref{prog:F12} gave you applies unchanged, because there is nothing
here for them to be different about.
\end{note}

\mermaidfig{p15-hold-the-rest-still}{Freezing the other inputs, and collecting
what is left.}{holding the rest still}
\end{fr}

\begin{fr}
Try it. $f(x, y) = x^{2}y + 3y$.

Differentiate with respect to $x$, treating $y$ as a constant.
\dotline
\nextframe
\end{fr}

\begin{fr}
\ans{$2xy$}
The first term is $y$ times $x^{2}$, and $y$ is a constant here, so it comes
out front: the derivative is $y \cdot 2x$. The second term, $3y$, has no $x$
in it at all, so as far as $x$ is concerned it is a constant and contributes
nothing.

Now the other one. $\frac{\partial f}{\partial y}$, holding $x$ still.
\nextframe
\end{fr}

\begin{fr}
\ans{$x^{2} + 3$}
This time $x^{2}$ is the constant multiplying $y$, so it survives as itself,
and $3y$ differentiates to $3$.

\begin{trapbox}
Look at what just happened to the two answers. $\frac{\partial f}{\partial x}$
still has a $y$ in it and $\frac{\partial f}{\partial y}$ still has an $x$.

\textbf{Holding an input still does not remove it from the answer.} It is
held still \emph{while you differentiate}, and then it goes back to being a
variable. A partial derivative is a function of all the same inputs the
original had \dash{} which is why it has to be evaluated at a point before it
is a number, and why \enquote{the derivative} on its own is not yet an
answer to anything.
\end{trapbox}
\end{fr}

\begin{fr}
So evaluate them. At the point $x = \val{p15.pt.x}$, $y = \val{p15.pt.y}$,
what are the two partial derivatives?
\dotline
\nextframe
\end{fr}

\begin{fr}
\ans{$\val{p15.fx.at}$ and $\val{p15.fy.at}$}
From $2xy = 2 \times \val{p15.pt.x} \times \val{p15.pt.y}$ and
$x^{2} + 3 = \val{p15.pt.x}^{2} + 3$.

Both were checked against a central difference over a grid of points, the way
Program~\ref{prog:F12} checked its rules, and they agree to better than
$\val{p15.partial.bound}$ everywhere on it. That is a check on the
arithmetic and not on the idea, but it is the check that catches a sign.

\begin{aibox}
This is what a framework computes for every parameter in a model. The loss is
one number; each weight is one input; and the partial derivative with respect
to one weight is the answer to \emph{if I nudge this one weight and nothing
else, how much does the loss move}.

There are as many of them as there are weights, and the interesting fact
about them is not any one number but that a single backward pass produces all
of them at once \dash{} which is Program~\ref{prog:P16}'s subject and needs
the shape of the computation rather than more calculus.
\end{aibox}
\end{fr}

% ============================================================
\section{Collecting them into one object}
\label{sec:P15-gradient}
\index{gradient}

\begin{fr}
Two partial derivatives at a point are two numbers. Write them as a vector,
in the order the inputs come:
\[ \nabla f = \left(\frac{\partial f}{\partial x},\
                    \frac{\partial f}{\partial y}\right) \]
That is the \textbf{gradient}, and $\nabla$ is read \enquote{del} or
\enquote{nabla}.

At the point from section 1 it is $(\val{p15.fx.at}, \val{p15.fy.at})$.
Program~\ref{prog:F09} gave a vector a length. What is this one's?
\dotline
\nextframe
\end{fr}

\begin{fr}
\ans{$\val{p15.grad.len}$}
From $\sqrt{\val{p15.fx.at}^{2} + \val{p15.fy.at}^{2}}$.

Hold on to the fact that it has a length as well as a direction, because both
are going to turn out to mean something, and the meanings are different.

\begin{notationbox}
\textbf{The word \enquote{gradient} is doing two jobs in this book and this is
where they meet.}

Program~\ref{prog:F06} and Program~\ref{prog:F11} use it in its ordinary
British sense: the \emph{slope} of a line, a single number, the $m$ in
$y = mx + c$. That was the right word there and those programs are not wrong.

Here it is a \emph{vector}, one component per input. The two are related
\dash{} in one dimension the gradient has one component and that component is
the slope \dash{} but they are not the same kind of object, and a sentence
like \enquote{the gradient is large} means something different depending on
which is meant. When this program says gradient, it means the vector.
\end{notationbox}
\end{fr}

\begin{fr}
A question that sounds trivial and is not. Why write two numbers as a vector
at all, rather than keeping them as two numbers?
\dotline
\nextframe
\end{fr}

\begin{fr}
\begin{ansblock}
Because a vector can be dotted with another vector, and two loose numbers
cannot. Everything in the rest of this program comes out of that one
operation.
\end{ansblock}

Program~\ref{prog:P05} spent a whole program on what a dot product buys, and
the answer was lengths, angles and perpendicularity from one definition.
Section 3 takes the gradient into that machinery and gets steepest descent
out without defining anything new.

\begin{note}
In more than two inputs nothing changes except the length of the list. A model
with a billion parameters has a gradient with a billion components, it is
still one vector, and every sentence in this program is still true of it.

The book stays in two dimensions because two is where you can see it. It is
worth remembering Program~\ref{prog:P05}'s warning while you do:
two-dimensional pictures of high-dimensional geometry mislead in specific
measurable ways, and the dot product is one of the few things that transfers
without a caveat.
\end{note}
\end{fr}

\begin{fr}
One more property before it is used. Program~\ref{prog:F11} said the
derivative of a constant is zero. What is the gradient at a point where $f$
has a minimum?
\nextframe
\end{fr}

\begin{fr}
\begin{ansblock}
The zero vector: every partial derivative is zero there.
\end{ansblock}

Because if any one of them were not, you could move a little way along that
input and go down, so you were not at a minimum.

\begin{trapbox}
The converse is not true, and Program~\ref{prog:F11} said so in one
dimension: a zero derivative is not a minimum. In more than one it gets
worse, because there is a shape that has no one-dimensional analogue at all
\dash{} a point that is a minimum along one direction and a maximum along
another.

Program~\ref{prog:P10} named it and drew it: a \emph{saddle}. Its gradient is
the zero vector, so anything that stops when the gradient vanishes stops
there quite happily. Telling the three apart takes second derivatives, which
is Program~\ref{prog:P17}'s.
\end{trapbox}
\end{fr}

% ============================================================
\section{The rate of change along any direction}
\label{sec:P15-directional}
\index{derivative!directional}

\begin{fr}
The two partial derivatives answer two questions: how fast $f$ changes if you
move along $x$, and if you move along $y$.

Neither is the question you actually have. You want to move in some direction
of your choosing \dash{} say north-east, or any of the infinitely many others
\dash{} and know what $f$ does.

Write a direction as a unit vector $u$. How would you define the rate of
change of $f$ along $u$, using only Program~\ref{prog:F11}'s definition of a
derivative?
\dotline
\nextframe
\end{fr}

\begin{fr}
\begin{ansblock}
Exactly as F11 did, along that line: take a small step $h$ in the direction
$u$, see how much $f$ changed, divide by $h$, and shrink $h$.
\end{ansblock}

\[ D_u f = \lim_{h \to 0} \frac{f(\vect{p} + h u) - f(\vect{p})}{h} \]

It is the same limit, on the one-variable function you get by walking along a
straight line through the point. Nothing new has been defined; a direction has
been chosen and F11's definition applied along it.

\textbf{The unit part matters.} If $u$ were twice as long the answer would
double, and it would be reporting the length of your ruler rather than the
steepness of the ground.
\end{fr}

\begin{fr}
Now the step this section exists for. That limit is a nuisance to compute for
every direction you might care about \dash{} and it turns out you never have
to.

The two partials are the answers for $u = (1, 0)$ and $u = (0, 1)$. A general
$u = (u_1, u_2)$ is $u_1$ of the first and $u_2$ of the second. So what is
$D_u f$, in terms of the gradient?
\dotline
\nextframe
\end{fr}

\begin{fr}
\begin{ansblock}
\[ D_u f = \nabla f \cdot u \]
\end{ansblock}

The rate along any direction is the gradient dotted with that direction, and
the two partials are the special cases $u = (1,0)$ and $u = (0,1)$.

That is the whole of this section, and it is worth being clear about what
kind of statement it is. It is not a definition and it is not a convention:
it is a consequence of the limit being linear in the direction, and it can be
checked. The script takes $\val{p15.dir.steps}$ directions right round the
circle, computes the limit numerically for each and compares it against the
dot product. The worst disagreement anywhere is under
$\val{p15.dir.bound}$.

\transcript{p15-directional}

\mermaidfig{p15-any-direction}{One number per direction, and where it comes
from.}{any direction}
\end{fr}

\begin{fr}
Before the consequences, one calculation. At the same point, with gradient
$(\val{p15.fx.at}, \val{p15.fy.at})$, what is the rate of change in the
direction $u = (\num{0.6}, \num{0.8})$?

Check that $u$ is a unit vector first, because if it is not the answer means
nothing.
\dotline
\nextframe
\end{fr}

\begin{fr}
\begin{ansblock}
$\val{p15.fx.at} \times \num{0.6} + \val{p15.fy.at} \times \num{0.8} =
\num{17.6}$, and $u$ is a unit vector because
$\num{0.6}^{2} + \num{0.8}^{2} = 1$.
\end{ansblock}

Which is the number the listing above prints, and it is smaller than the
gradient's own length of $\val{p15.grad.len}$. That is not a coincidence and
the next section is about why.

\begin{note}
The $3$--$4$--$5$ triangle is why $(\num{0.6}, \num{0.8})$ is a unit vector,
and Program~\ref{prog:P09} and Program~\ref{prog:P10} both used the same
triple to build rotations with no rounding in them. It is the cheapest exact
unit vector there is, which is why it keeps appearing.
\end{note}
\end{fr}

% ============================================================
\section{Steepest, and across the contour}
\label{sec:P15-steepest}
\index{gradient!steepest direction}

\begin{fr}
Now read $D_u f = \nabla f \cdot u$ with Program~\ref{prog:P05}'s identity
in front of you:
\[ a \cdot b = \lVert a \rVert \, \lVert b \rVert \cos\theta \]
Since $u$ is a unit vector, $\lVert u \rVert = 1$, so
\[ D_u f = \lVert \nabla f \rVert \cos\theta \]
where $\theta$ is the angle between the direction you chose and the gradient.

Look at what that has done to the question. The first factor is the same
whatever direction you pick \dash{} it is a property of the point, not of your
choice \dash{} so the only thing your choice controls is $\cos\theta$. The
whole of \enquote{which of the infinitely many directions} has collapsed into
\enquote{which cosine}, and that is a question with a one-word answer.

Which $\theta$ makes the rate biggest, and what is the biggest rate?
\dotline
\nextframe
\end{fr}

\begin{fr}
\begin{ansblock}
$\theta = 0$, and the biggest rate is $\lVert \nabla f \rVert$ itself.
\end{ansblock}

Because $\cos\theta$ is at most $1$ and equals $1$ only when the two point
the same way.

\textbf{So the gradient is the direction of steepest increase, and this is
derived rather than asserted.} It is not a definition and it is not a fact
you have to take on trust: it follows from the rate being a dot product, and
from the cosine being largest at zero. That is the sentence this whole
program exists to earn.

And the length now means something too: $\lVert \nabla f \rVert$ is
\emph{how fast} $f$ climbs in that best direction. Direction and length,
two meanings, one vector.
\end{fr}

\begin{fr}
The measurement, because \enquote{the biggest} is a claim about every
direction rather than about one. Over the $\val{p15.dir.steps}$ directions,
the largest rate found is the one at $\val{p15.grad.deg}$ degrees, which is
the gradient's own direction, and its value is the gradient's own length.

The sampled maximum falls short of $\lVert \nabla f \rVert$ by under
$\val{p15.sweep.shortfall.bound}$, and the script asserts that the shortfall
is no more than the grid spacing can account for \dash{} not that it is zero,
which would be false for any finite sweep.

\begin{note}
The first version of that check asserted the sampled maximum \emph{equalled}
the length, and it failed. It deserved to: a maximum over three and a half
thousand directions cannot equal a maximum over all of them, so the tolerance
was a number chosen to make a check pass rather than a statement about the
mathematics.

What is asserted now is a pair \dash{} that no direction beats the gradient,
which is exact and is Program~\ref{prog:P05}'s inequality; and that the
sampled best falls short by no more than the resolution. The second turns the
shortfall from an embarrassment into an explanation.
\end{note}
\end{fr}

\begin{fr}
Now the other end of the same equation. Which $\theta$ makes the rate zero,
and what does that mean about where you can walk?
\dotline
\nextframe
\end{fr}

\begin{fr}
\begin{ansblock}
$\theta = 90$ degrees. Walk at right angles to the gradient and $f$ does not
change at all, to first order.
\end{ansblock}

A path along which $f$ does not change is a \textbf{contour} \dash{} a level
set, the lines on a map, the rings round the bottom of a bowl. So:

\textbf{the gradient is perpendicular to the contour through the point.}

That is the same equation read at its other extreme, and it is worth noticing
that \emph{steepest uphill} and \emph{perpendicular to the contour} are not
two facts about the gradient. They are one fact, read at $\cos\theta = 1$ and
at $\cos\theta = 0$.
\end{fr}

\begin{fr}
That claim can be made exactly rather than numerically, and the script does.
On $f = ax^{2} + by^{2}$ the gradient is $(2ax, 2by)$, and a direction along
the contour is $(-by, ax)$. Dot them together.
\nextframe
\end{fr}

\begin{fr}
\begin{ansblock}
$-2abxy + 2abxy = 0$. Not approximately zero \dash{} identically zero, for
every $a$, $b$, $x$ and $y$.
\end{ansblock}

The script checks it over the rationals at $\val{p15.perp.trials}$
combinations, with no tolerance anywhere, because a claim containing the word
\emph{perpendicular} is a claim about an exact zero and a floating-point
comparison would have made it a claim about a threshold.

\begin{note}
That is the same discipline Program~\ref{prog:P04} used for rank and
Program~\ref{prog:P13} used for a stationary distribution: where the
mathematics is exact, compute it exactly and assert equality. A tolerance is
for a measurement. This is not one.
\end{note}
\end{fr}

\begin{fr}
Here is why that is about to matter, and it is worth answering before reading
on.

You are standing on a bowl-shaped surface, somewhere off to one side. The
gradient at your feet points at right angles to the contour you are standing
on. Does the direction straight down that slope point at the lowest point of
the bowl?
\dotline
\nextframe
\end{fr}

\begin{fr}
\begin{ansblock}
Only if the contours are circles. In general, no.
\end{ansblock}

Most people answer yes, and the reason is a good one: it is true for the bowl
everyone pictures. A round bowl has circular contours, the perpendicular to a
circle points at its centre, and the two directions coincide exactly.

\textbf{But perpendicular-to-the-contour is a local statement and
where-the-minimum-is is a global one}, and nothing makes them agree. Stretch
the bowl along one axis and the contours become ellipses; the perpendicular to
an ellipse does not point at its centre unless you happen to be standing on
one of its axes.

Section 5 is what that does to a walk.

\mermaidfig{p15-across-the-valley}{Where the perpendicular points, in a round
bowl and in a long one.}{across the valley}
\end{fr}

% ============================================================
\section{What it does to the walk}
\label{sec:P15-zigzag}
\index{gradient!descent}

\begin{fr}
Program~\ref{prog:F11} already wrote the walk, in one dimension:
\[ x \;\longleftarrow\; x - \eta\, f'(x) \]
and called the minus sign the entire steering mechanism. In more than one
dimension the only thing that changes is that the derivative is a vector:
\[ \vect{p} \;\longleftarrow\; \vect{p} - \eta\, \nabla f(\vect{p}) \]

Notice that the update is \emph{componentwise}: each coordinate moves by its
own partial derivative and consults none of the others. That is why a
framework can apply this to a billion parameters without anything having to
be coordinated \dash{} the vector notation is a convenience for writing it
down, and the arithmetic underneath is a billion independent subtractions.

Say why the minus sign is there, in terms of this section's own result rather
than F11's.
\dotline
\nextframe
\end{fr}

\begin{fr}
\begin{ansblock}
Because the gradient is the direction of steepest \emph{increase}, and you
want to go down.
\end{ansblock}

Which is a stronger statement than F11 could make. F11 knew that against the
sign of the derivative is downhill; it did not know that the negative gradient
is the \emph{best} downhill direction available, because that needs the dot
product and the cosine. Now it is not a rule of thumb but the answer to
\emph{which of all the directions goes down fastest}.

\begin{trapbox}
\textbf{\enquote{Gradient descent moves in the direction of the gradient.}}
It moves in the direction of the \emph{negative} gradient, and this is the
single most common sign error in a hand-written training loop.

It is worth seeing what it costs rather than being told. At the point
$(\val{p15.zig.start.x}, \val{p15.zig.start.y})$ on the bowl below, the value
is $\val{p15.sign.before}$. One small step against the gradient takes it to
$\val{p15.sign.down}$. One small step \emph{with} the gradient takes it to
$\val{p15.sign.up}$.

Nothing raises an error. The loop runs, the numbers are finite, and the loss
climbs \dash{} which reads as a bad learning rate or a bad initialisation far
more often than it reads as a missing minus sign.
\end{trapbox}
\end{fr}

\begin{fr}
Now the bowl. Program~\ref{prog:P10} built one and deliberately did not spend
it: a quadratic whose two eigenvalues are $\val{p15.lam.hi}$ and
$\val{p15.lam.lo}$, so its contours are ellipses
$\val{p10.basin.axisratio}$ times longer than they are wide. That is a
valley: steep across, shallow along.

This program uses that same bowl rather than inventing another, and the script
reads P10's committed numbers and fails the build if they have moved.

Start at $(\val{p15.zig.start.x}, \val{p15.zig.start.y})$, which is off both
axes. Before computing anything: does the negative gradient there point at the
minimum?
\nextframe
\end{fr}

\begin{fr}
\ans{No \dash{} it is $\val{p15.zig.angle}$ degrees away from it}
The minimum is at the origin, so the direction home is straight back along the
line you are on. The negative gradient is perpendicular to the contour
instead, and on an ellipse those are two different directions.

\textbf{Forty-two degrees is not a rounding error, it is nearly half a right
angle.} The steering is not slightly off; it is pointing somewhere else
entirely, and it is doing so for a reason that is now completely understood
rather than mysterious.

\begin{note}
The round bowl is the control, and it is exact rather than nearly so. Set the
two eigenvalues equal and the gradient becomes a multiple of the position
itself, so the negative gradient points \emph{exactly} at the minimum and the
walk is a dead straight line. The script asserts that with a cross product
over floats, which is exactly zero.

Asking the arccosine for the angle does not say so: it reports about
$\val{p15.round.acos.bound}$ of a degree rather than zero, because the cosine
comes back as one minus a rounding and the arccosine of a number that close to
$1$ is where Program~\ref{prog:P02}'s cancellation lives. The geometry is
exact; the arccosine is the wrong instrument for measuring a zero.
\end{note}
\end{fr}

\begin{fr}
So take the walk. $\val{p15.zig.steps}$ steps of size $\val{p15.zig.eta}$
from that corner, on that bowl.

Two things to predict before reading on. How many times does the walk cross
from one side of the valley to the other? And how does the distance it travels
compare with the distance it actually covers?
\dotline
\nextframe
\end{fr}

\begin{fr}
\begin{ansblock}
It crosses $\val{p15.zig.crossings}$ times in $\val{p15.zig.steps}$ steps
\dash{} every single step \dash{} and travels $\val{p15.zig.ratio}$ times as
far as it moves.
\end{ansblock}

$\val{p15.zig.pathlen}$ of path to cover $\val{p15.zig.moved}$ of ground.
Nearly seven-eighths of the walking is sideways.

\begin{note}
The measure is distance travelled against distance \emph{covered}, not against
the distance to the minimum. An earlier draft used the second and reported a
figure below one for a walk that is a dead straight line \dash{} because after
twenty steps it has not arrived, so it has travelled less than the whole way
home. That says something about the step size and nothing about the direction.

Path over displacement is $1$ for any straight path however far it gets, which
is exactly the property this section is about, and it is $1$ for the round
bowl to the last bit.
\end{note}
\end{fr}

\begin{fr}
The mechanism is one line of arithmetic and it is Program~\ref{prog:F11}'s own
recurrence, one direction at a time.

On this bowl the two directions do not interact: each coordinate is multiplied
by $1 - \eta\lambda$ every step, with its own $\lambda$. With
$\eta = \val{p15.zig.eta}$, work out that factor for the steep direction
($\lambda = \val{p15.lam.hi}$) and for the shallow one
($\lambda = \val{p15.lam.lo}$).
\dotline
\nextframe
\end{fr}

\begin{fr}
\ans{$\val{p15.fac.hi}$ and $\val{p15.fac.lo}$}
And there is the zig-zag, in the signs. The steep coordinate is multiplied by
a \emph{negative} number, so it changes side every single step while shrinking
by a fifth. The shallow one is multiplied by a positive number just under one,
so it creeps towards zero without ever crossing.

\textbf{One direction overshoots and one crawls, and the same step size is
responsible for both.} Make the step small enough to stop the overshooting and
the shallow direction takes forever; make it big enough to move the shallow
direction and the steep one flies apart.

\begin{rigourbox}
What the largest usable $\eta$ actually is, and how it depends on the
curvature, is Program~\ref{prog:P17}'s: it needs the second derivative
properly, which is a matrix and has not been defined. This program can say
that the factor goes negative and that the walk crosses; it may not tell you
where the boundary is.

And the fix is Program~\ref{prog:P20}'s. Momentum averages the recent steps,
so the sideways components \dash{} which alternate in sign \dash{} cancel
against each other while the forward ones \dash{} which do not \dash{} add up.
That sentence is only obvious once you have seen this picture, which is why
this program produces the problem and leaves it standing.
\end{rigourbox}
\end{fr}

\begin{fr}
One last thing, because it is the reason any of this generalises.

Everything in this section was worked on two inputs and a bowl you can draw. A
real loss has a billion inputs and cannot be drawn at all. Which of these
sentences survives that, and which was about the picture?
\dotline
\nextframe
\end{fr}

\begin{fr}
\begin{ansblock}
All of them survive. Not one step used the fact that there were two.
\end{ansblock}

The directional derivative is a dot product in any number of dimensions; the
cosine is largest at zero in any number; a contour becomes a surface and the
gradient is still perpendicular to it; and a valley becomes a direction in
which the curvature is small compared with another, of which there are
plentifully many when there are a billion to choose from.

\textbf{What does not survive is the picture, and it was never the argument.}
The drawing was a way of seeing a dot product. The dot product is what is
true.

\begin{aibox}
This is why \enquote{the loss landscape} is a useful phrase and a dangerous
one. The steering really is what this program described, in as many dimensions
as you like. The \emph{shape} is not: Program~\ref{prog:P05} measured how
badly two-dimensional intuition transfers, and a billion-dimensional surface
has room for behaviour that has no two-dimensional counterpart at all.

Trust the algebra here and distrust the mental image, in that order.
\end{aibox}
\end{fr}

% ============================================================
\begin{summarybox}
\sumitem{1--2}{\result{\textbf{A partial derivative holds every other input still}
  and is then an ordinary one-variable derivative, so every rule from
  Program~\ref{prog:F12} applies unchanged}.}
\sumitem{3--5}{\result{For $f = x^{2}y + 3y$ the partials are $2xy$ and $x^{2}+3$.
  \textbf{Holding an input still does not remove it from the answer} \dash{}
  a partial derivative is still a function of everything}.}
\sumitem{6--7}{\result{At $(\val{p15.pt.x}, \val{p15.pt.y})$ they are
  $\val{p15.fx.at}$ and $\val{p15.fy.at}$, checked against a central
  difference over a grid to better than $\val{p15.partial.bound}$}.}
\sumitem{8--9}{\result{\textbf{The gradient is the partials written as a vector},
  $\nabla f$, with a length as well as a direction \dash{} here
  $\val{p15.grad.len}$}.}
\sumitem{9}{\result{\textbf{The word does two jobs.} Programs F06 and F11 use
  \enquote{gradient} for the slope of a line, a single number; here it is a
  vector with one component per input}.}
\sumitem{10--11}{\result{It is written as a vector \textbf{so that it can be dotted
  with one}, which is where the rest of the program comes from}.}
\sumitem{12--13}{\result{At a minimum the gradient is the zero vector, and
  \textbf{the converse fails} \dash{} a saddle has a zero gradient too, and
  telling them apart is Program~\ref{prog:P17}'s}.}
\sumitem{14--17}{\result{\textbf{The rate of change along a unit direction $u$ is
  $\nabla f \cdot u$}, which is Program~\ref{prog:F11}'s limit taken along a
  line. Checked against that limit over $\val{p15.dir.steps}$ directions to
  better than $\val{p15.dir.bound}$}.}
\sumitem{20--21}{\result{With Program~\ref{prog:P05}'s cosine,
  $D_u f = \lVert\nabla f\rVert\cos\theta$, so \textbf{the steepest direction
  is the gradient's own and the fastest rate is its length}. Derived, not
  asserted}.}
\sumitem{23--24}{\result{At $\cos\theta = 0$ the rate is zero, so \textbf{the
  gradient is perpendicular to the contour} \dash{} the same equation read at
  its other end, exactly zero over $\val{p15.perp.trials}$ rational cases}.}
\sumitem{25--26}{The perpendicularity is checked \textbf{exactly over the
  rationals}, with no tolerance anywhere, because a claim containing the word
  \emph{perpendicular} is a claim about an exact zero.}
\sumitem{27--28}{\result{\textbf{Perpendicular to the contour is not the same as
  pointing at the minimum}, unless the contours are circles. One is local and
  one is global}.}
\sumitem{29--30}{\result{The walk is Program~\ref{prog:F11}'s with a vector
  derivative, and the minus sign now has a stronger reason: the gradient is
  steepest \emph{increase}}.}
\sumitem{30}{\result{\textbf{The sign is the commonest error in a hand-written
  loop.} One step the right way takes $\val{p15.sign.before}$ to
  $\val{p15.sign.down}$; one step the wrong way takes it to
  $\val{p15.sign.up}$, with no error raised}.}
\sumitem{31--32}{\result{On Program~\ref{prog:P10}'s bowl, eigenvalues
  $\val{p15.lam.hi}$ and $\val{p15.lam.lo}$, the negative gradient points
  $\val{p15.zig.angle}$ degrees away from the minimum. On a round bowl it
  points exactly at it}.}
\sumitem{33--34}{\result{The walk crosses the valley $\val{p15.zig.crossings}$ times
  in $\val{p15.zig.steps}$ steps and \textbf{travels $\val{p15.zig.ratio}$
  times as far as it moves}}.}
\sumitem{35--36}{\result{Each coordinate is multiplied by $1 - \eta\lambda$: the
  factor is $\val{p15.fac.hi}$ in the steep direction, which flips its sign
  every step, and $\val{p15.fac.lo}$ in the shallow one, which creeps.
  \textbf{One step size has to serve both.}}}
\sumitem{37--38}{\result{\textbf{Not one step of the argument used the fact that
  there were two inputs.} The picture is an illustration of the dot product
  and never a premise of it}.}
\end{summarybox}

\canyou

\begin{testexercises}
\item $g(x, y) = 3x^{2}y - y^{3}$. Write both partial derivatives, and
  evaluate the gradient at $(1, 2)$.
  \answerto{$\partial g/\partial x = 6xy$ and $\partial g/\partial y =
  3x^{2} - 3y^{2}$. At $(1, 2)$ that is
  \[ \nabla g = (12,\ -9), \qquad \lVert \nabla g \rVert = 15. \]
  A $3$--$4$--$5$ triangle again, scaled by three.}
\item At a point where $\nabla f = (3, 4)$, what is the rate of change of $f$
  in the direction $(\num{0.8}, \num{0.6})$, and in the direction
  $(-\num{0.8}, -\num{0.6})$?
  \answerto{$\num{4.8}$ and $-\num{4.8}$. Reversing the direction reverses the
  sign and nothing else, because the dot product is linear in $u$ \dash{} which
  is also why the steepest \emph{descent} direction is exactly opposite the
  steepest ascent one rather than merely near it.}
\item A colleague says their loss is falling but very slowly, and that the
  gradient's length is large. What can you tell them from this program alone?
  \answerto{That large gradient and slow progress are not a contradiction, and
  the likely reason is that most of the gradient is pointing across a valley
  rather than along it. The length says how steep the ground is in the
  steepest direction; it says nothing about whether that direction goes
  anywhere useful. What to do about it is Program~\ref{prog:P20}'s.}
\item Why is it wrong to compare $D_u f$ for two directions if one of them is
  not a unit vector?
  \answerto{Because $D_u f$ scales with the length of $u$, so a longer $u$
  reports a bigger number without the ground being any steeper. You would be
  comparing rulers rather than slopes. The unit requirement is what makes the
  number a property of the surface.}
\item On a bowl with contours that are circles, how many times does gradient
  descent cross from one side to the other before arriving?
  \answerto{None, for any step size that converges at all. The negative
  gradient points exactly at the centre, so every step is along the same
  straight line and the path is $1$ times its own displacement \dash{} which
  the script asserts exactly. Crossing needs the contours to be non-circular,
  and that is the whole content of section 5.}
\end{testexercises}

\begin{furtherproblems}
\item Show that if $u$ is a unit vector then $D_{-u} f = -D_u f$, and say what
  that means about the steepest descent direction.
  \answerto{$\nabla f \cdot (-u) = -(\nabla f \cdot u)$, because the dot
  product is linear in its second argument. So the direction of steepest
  \emph{decrease} is exactly $-\nabla f$: the same argument that made
  $\cos\theta = 1$ the best case makes $\cos\theta = -1$ the worst, and the
  worst rate of increase is the best rate of decrease.}
\item The gradient of $f(x,y) = ax^{2} + by^{2}$ at $(x, y)$ is $(2ax, 2by)$.
  For which points is it parallel to the direction back to the origin?
  \answerto{Where $2ax \cdot y = 2by \cdot x$, that is $xy(a - b) = 0$. So
  either $a = b$ \dash{} a round bowl, where it is parallel everywhere
  \dash{} or the point lies on one of the two axes. Everywhere else on an
  elongated bowl the two directions differ, which is why a walk that starts
  off the axes zig-zags and a walk that starts on one does not.}
\item A function of three inputs has a gradient with three components. How
  many directions are at right angles to it, and what does that mean about
  the contour?
  \answerto{A whole plane of them rather than two isolated directions: in
  three dimensions the vectors perpendicular to a given one form a
  two-dimensional space, which is Program~\ref{prog:P04}'s dimension counting.
  So the level set through a point is a \emph{surface} rather than a curve,
  and the gradient is perpendicular to all of it at once.}
\item Why does the argument in section 4 never mention how many inputs $f$
  has?
  \answerto{Because it uses only that $D_u f$ is a dot product and that a dot
  product is $\lVert a\rVert\lVert b\rVert\cos\theta$, and neither statement
  has a dimension in it. That is what makes the result usable on a model with
  a billion parameters, where nothing can be drawn. The picture in section 5
  is an illustration of the algebra and never a premise of it.}
\item You have a bowl whose eigenvalues are $\lambda$ and $1$, and you double
  $\lambda$. What happens to the angle between the negative gradient and the
  direction home, and how would you check?
  \answerto{It grows: the more elongated the bowl, the further the
  perpendicular points from the centre. The script checks exactly that, by
  sweeping the ratio and asserting the angle increases at every step \dash{}
  which is the invariant, where any single angle would be a fact about one
  bowl. It cannot grow past a right angle, because the step would then be
  going uphill.}
\end{furtherproblems}
